\documentclass[a4paper,12pt]{report}
\PassOptionsToPackage{hyphens}{url}
\usepackage{graphicx}
\usepackage{listings}
\usepackage{caption}
\usepackage{amsmath}
\usepackage{fancyhdr}
\usepackage[hidelinks]{hyperref}
\usepackage{url}
\usepackage{titlesec}
\usepackage[a4paper, margin=1in]{geometry}
\usepackage{float}
\DeclareUnicodeCharacter{2212}{-}
\usepackage{makecell}
\renewcommand\cellalign{lc}
\usepackage{booktabs}

\setlength{\parindent}{0pt}
\lstset{
basicstyle=\ttfamily\footnotesize,
columns=fullflexible,
keepspaces=true,
breaklines=true,
breakatwhitespace=true,
postbreak=\mbox{\textcolor{red}{$\hookrightarrow$}\space},
}

\hypersetup{
colorlinks=false,
pdfborder={0 0 0},
linkcolor=black,
filecolor=black,
urlcolor=black
}

\titleformat{\chapter}[hang]{\normalfont\huge\bfseries}{\thechapter}{2pc}{}
\titlespacing*{\chapter}{0pt}{-60pt}{20pt}

\title{Real-Time Drone Detection Using YOLOv26n\\
\vspace{0.5cm}\href{https://github.com/AEEltayeb/drone_detection_py}{\texttt{GitHub Repository}}}
\author{Ahmed Eltayeb}
\date{\today}

\begin{document}

\maketitle

% Abstract
\begin{abstract}
This project presents a real-time drone detection system based on YOLOv26n, a state-of-the-art object detection model. The system is trained on a carefully curated dataset of 172,022 images from 10 different sources, combining drone imagery with background datasets to suppress false positives. The trained model achieves a precision of 97.8\%, recall of 91.5\%, and mAP50 of 95.1\% on the test set, with inference speeds of approximately 2.2ms per image. The project includes a desktop GUI application for real-time detection demonstration and ESP32 integration for physical alert notification via buzzer and OLED display.
\end{abstract}

% Table of contents
\tableofcontents
\newpage

% List of Figures
\listoffigures
\newpage

% List of Tables
\listoftables
\newpage

% Introduction
\chapter{Introduction}

Unmanned Aerial Vehicles (UAVs), commonly known as drones, have become increasingly prevalent in both civilian and commercial applications. However, the widespread availability of drones has also raised significant security concerns, including privacy violations, smuggling, and potential threats to critical infrastructure. This has created an urgent need for reliable and efficient drone detection systems.

This project addresses the challenge of real-time drone detection using deep learning-based computer vision. Specifically, we employ YOLOv26n (You Only Look Once), a compact and efficient variant of the YOLO family of object detectors, fine-tuned on a comprehensive dataset of drone imagery.

The key challenges in drone detection include:
\begin{itemize}
    \item \textbf{Small object detection:} Drones often appear as small objects in surveillance footage, making detection difficult.
    \item \textbf{False positive suppression:} Birds and other flying objects can be easily confused with drones.
    \item \textbf{Real-time performance:} Detection must be fast enough for practical deployment.
    \item \textbf{Varying conditions:} Drones must be detected across different lighting conditions, backgrounds, and distances.
\end{itemize}

Our approach combines careful dataset curation, transfer learning from pretrained YOLO weights, and a real-time detection pipeline with temporal filtering to address these challenges. The system also integrates with an ESP32 microcontroller to provide physical alerts when drones are detected.

\newpage

% Objectives
\chapter{Objectives}

\section{Project Pipeline}

The objectives of this project are organized as a sequential pipeline to ensure systematic development of a robust drone detection system:

\begin{enumerate}
    \item \textbf{Dataset Collection and Merging:} Collect drone imagery from multiple public datasets, including video-based and image-based sources. Merge positive samples (drones) with negative samples (birds, vehicles, backgrounds) to create a balanced training set.
    
    \item \textbf{Data Preprocessing:} Extract frames from video datasets at appropriate sampling rates, convert annotations to YOLO format, and resize images to the target resolution (640$\times$640).
    
    \item \textbf{Model Fine-tuning:} Fine-tune the YOLOv26n model using pretrained weights, optimizing hyperparameters for drone detection performance.
    
    \item \textbf{Real-Time Detection System:} Develop a detection pipeline with temporal filtering, cascaded ROI confirmation, and warning/alert logic to reduce false positives.
    
    \item \textbf{Hardware Integration:} Integrate with ESP32 microcontroller for physical alert notification via buzzer and OLED display.
    
    \item \textbf{GUI Application:} Develop a desktop GUI for demonstration and configuration of the detection system.
\end{enumerate}

\newpage

% Dataset Preparation
\chapter{Dataset Preparation}

\section{Dataset Sources}

For this project, we compiled a comprehensive dataset from 10 different sources, totaling \textbf{172,022 images}. The datasets are divided into two categories: drone datasets (positive samples) and false positive suppression datasets (negative samples).

\subsection{Drone Datasets (Positive Samples)}

\begin{table}[H]
\centering
\caption{Drone Dataset Sources}
\label{tab:drone_datasets}
\begin{tabular}{|l|l|l|l|}
\hline
\textbf{Dataset} & \textbf{Type} & \textbf{Sampling} & \textbf{Notes} \\
\hline
wosdetc & Video & 1 in 2 frames & Bird vs Drone; birds used as background \\
\hline
Anti-UAV-RGBT & Video & 1 in 5 frames & Visible only; infrared not used \\
\hline
DUT Anti-UAV & Images & All & Image dataset \\
\hline
Anti-MUAV-roboflow & Images & All & Image dataset \\
\hline
mav-vid & Images & All & Image dataset \\
\hline
\end{tabular}
\end{table}

\subsection{False Positive Suppression Datasets}

To reduce false positives, particularly from birds and other common objects, we included the following background datasets:

\begin{table}[H]
\centering
\caption{False Positive Suppression Datasets}
\label{tab:fp_datasets}
\begin{tabular}{|l|l|l|}
\hline
\textbf{Dataset} & \textbf{Type} & \textbf{Purpose} \\
\hline
airbird & Images & Bird-only dataset to suppress bird false positives \\
\hline
bdd10k & Images & IoV (Internet of Vehicles) dataset \\
\hline
fbd-sv & Video (1 in 3 frames) & Bird-only dataset \\
\hline
DETRAC & Images & CCTV footage of streets and junctions \\
\hline
virat & Video (1 in 3 frames) & Surveillance video dataset \\
\hline
\end{tabular}
\end{table}

\section{Data Processing}

All video datasets were processed to extract frames at the specified sampling rates. The extracted frames were then:
\begin{itemize}
    \item Resized to 640$\times$640 pixels (YOLO input size)
    \item Converted to YOLO annotation format (normalized bounding box coordinates)
    \item Split into training and validation sets
\end{itemize}

\begin{figure}[H]
\centering
\includegraphics[width=0.9\textwidth]{images/labels.jpg}
\caption{Dataset Label Distribution and Bounding Box Statistics}
\label{fig:labels}
\end{figure}

\newpage

% Model Training
\chapter{Model Training}

\section{Training Configuration}

The model was trained using YOLOv26n with pretrained weights, fine-tuned on our drone detection dataset. Table~\ref{tab:training_hyperparams} summarizes the key hyperparameters used.

\begin{table}[H]
\centering
\caption{Training Hyperparameters}
\label{tab:training_hyperparams}
\begin{tabular}{|l|c|}
\hline
\textbf{Hyperparameter} & \textbf{Value} \\
\hline
Model Architecture & YOLOv26n (pretrained) \\
Epochs & 70 \\
Batch Size & 48 \\
Image Size & 640$\times$640 \\
Learning Rate (lr0) & 0.001 \\
Optimizer & AdamW \\
Mosaic Augmentation & 0.33 \\
Close Mosaic & 20 \\
\hline
\end{tabular}
\end{table}

\subsection{Hardware Configuration}
\begin{itemize}
    \item \textbf{GPU:} NVIDIA RTX 3090 (cloud rented)
    \item \textbf{Training Time:} $\sim$18 hours
    \item \textbf{Best Epoch:} 68
\end{itemize}

\section{Training Results}

\subsection{Training Curves}

Figure~\ref{fig:training_curves} shows the training and validation metrics across 70 epochs. Both box loss and classification loss decrease steadily, with validation metrics closely tracking training metrics, indicating no significant overfitting.

\begin{figure}[H]
\centering
\includegraphics[width=0.95\textwidth]{images/reconstructed_results.png}
\caption{Training and Validation Metrics Across Epochs}
\label{fig:training_curves}
\end{figure}

\subsection{Quantitative Results}

Table~\ref{tab:train_metrics} shows the evolution of key metrics across selected epochs.

\begin{table}[H]
\centering
\caption{Training Metrics at Selected Epochs}
\label{tab:train_metrics}
\begin{tabular}{|r|r|r|r|r|r|r|}
\hline
\textbf{Epoch} & \textbf{Train Box} & \textbf{Train Cls} & \textbf{Precision} & \textbf{Recall} & \textbf{mAP50} & \textbf{mAP50-95} \\
\hline
1 & 2.013 & 1.373 & 0.788 & 0.652 & 0.745 & 0.384 \\
10 & 1.673 & 0.700 & 0.940 & 0.867 & 0.911 & 0.552 \\
30 & 1.497 & 0.545 & 0.969 & 0.905 & 0.941 & 0.609 \\
50 & 1.382 & 0.486 & 0.974 & 0.918 & 0.950 & 0.638 \\
68 & 1.256 & 0.421 & 0.979 & 0.920 & 0.955 & 0.653 \\
70 & 1.248 & 0.416 & 0.978 & 0.921 & 0.955 & 0.654 \\
\hline
\end{tabular}
\end{table}

\subsection{Final Validation Results}

The final model (epoch 68) was evaluated on the test set with the following results:

\begin{table}[H]
\centering
\caption{Final Test Set Performance}
\label{tab:final_results}
\begin{tabular}{|l|c|}
\hline
\textbf{Metric} & \textbf{Value} \\
\hline
Precision (P) & 0.978 \\
Recall (R) & 0.915 \\
mAP50 & 0.951 \\
\hline
\end{tabular}
\end{table}

\subsection{Inference Speed}

\begin{table}[H]
\centering
\caption{Inference Speed Breakdown}
\label{tab:inference_speed}
\begin{tabular}{|l|c|}
\hline
\textbf{Stage} & \textbf{Time (ms)} \\
\hline
Preprocess & 0.7 \\
Inference & 1.4 \\
Loss & 0.0 \\
Postprocess & 0.1 \\
\hline
\textbf{Total} & \textbf{2.2} \\
\hline
\end{tabular}
\end{table}

\textbf{Effective FPS:} 25--35 FPS depending on image resolution (max 3 targets per frame).

\subsection{Confusion Matrix}

\begin{figure}[H]
\centering
\includegraphics[width=0.7\textwidth]{images/confusion_matrix_normalized.png}
\caption{Normalized Confusion Matrix}
\label{fig:confusion_matrix}
\end{figure}

\subsection{Precision-Recall Curve}

\begin{figure}[H]
\centering
\includegraphics[width=0.7\textwidth]{images/BoxPR_curve.png}
\caption{Precision-Recall Curve}
\label{fig:pr_curve}
\end{figure}

\subsection{Sample Detection Results}

\begin{figure}[H]
\centering
\includegraphics[width=0.95\textwidth]{images/val_batch0_pred.jpg}
\caption{Sample Detection Results on Validation Set}
\label{fig:sample_detections}
\end{figure}

\newpage

% Real-Time Detection System
\chapter{Real-Time Detection System}

\section{Detection Pipeline}

The real-time detection system implements several mechanisms to improve detection reliability and reduce false positives:

\subsection{Frame-Based Inference}
The system performs YOLO inference at a configurable frame rate (default: 5 FPS inference stride). Between inference frames, the previous detection results are displayed to maintain visual continuity.

\subsection{Temporal ROI Propagation}
When a drone is detected in a frame, the system creates a Region of Interest (ROI) around the detection. In subsequent frames, inference is performed on this ROI region first, improving detection consistency for moving targets.

\subsection{Cascaded ROI Confirmation}
For low-confidence or small detections, the system performs a second inference pass on a cropped ROI at higher resolution to confirm the detection before triggering an alert.

\subsection{Warning and Alert Logic}
The system uses a sliding window approach:
\begin{itemize}
    \item \textbf{Warning:} Triggered when 9 out of 10 inference frames contain a detection
    \item \textbf{Alert:} Triggered when the detection meets size and confidence thresholds
    \item \textbf{Cooldown:} 3-second cooldown between alerts to prevent spam
\end{itemize}

\section{ESP32 Integration}

The system integrates with an ESP32 microcontroller for physical alert notification:

\begin{itemize}
    \item \textbf{HTTP Server:} ESP32 runs a web server on port 5000
    \item \textbf{Endpoints:}
    \begin{itemize}
        \item \texttt{POST /alert/drone} -- Trigger drone alert (buzzer + display)
        \item \texttt{POST /alert/warning} -- Trigger warning (display only)
        \item \texttt{POST /alert/clear} -- Clear active alerts
        \item \texttt{GET /status} -- Get current system status
    \end{itemize}
    \item \textbf{Hardware:}
    \begin{itemize}
        \item OLED display (SSD1306, 128$\times$64)
        \item Passive buzzer (2kHz tone)
        \item LED indicator
        \item Dismiss button
    \end{itemize}
\end{itemize}

\section{GUI Application}

A desktop GUI application was developed using PySide6 (Qt) for demonstration and configuration:

\begin{itemize}
    \item Real-time video display via OpenCV
    \item Configuration panel for detection parameters
    \item ESP32 connection management
    \item Event log display
\end{itemize}

The GUI allows users to:
\begin{itemize}
    \item Select video source
    \item Configure detection thresholds
    \item Enable/disable temporal filtering
    \item Test ESP32 connection
    \item Save detection results
\end{itemize}

\newpage

% Conclusion
\chapter{Conclusion}

This project successfully developed a real-time drone detection system using YOLOv26n. The key achievements are:

\begin{itemize}
    \item Curated a comprehensive dataset of 172,022 images from 10 sources
    \item Achieved 97.8\% precision and 91.5\% recall on the test set
    \item Implemented real-time detection at 25--35 FPS
    \item Developed a complete system with GUI and hardware integration
\end{itemize}

\section{Limitations}
\begin{itemize}
    \item Detection performance degrades for very distant or occluded drones
    \item The system requires good lighting conditions for optimal performance
    \item Training was limited to available GPU resources
\end{itemize}

\section{Future Work}
\begin{itemize}
    \item Extend to multi-class detection (different drone types)
    \item Implement object tracking for continuous monitoring
    \item Deploy on edge devices (Jetson, Raspberry Pi)
    \item Explore thermal/infrared fusion for night detection
\end{itemize}

\newpage

% References
\chapter{References}

\section{Dataset Sources}
\begin{itemize}
    \item wosdetc Challenge: \url{https://github.com/wosdetc/challenge}
    \item Anti-UAV Dataset: \url{https://github.com/ZhaoJ9014/Anti-UAV}
    \item DUT Anti-UAV: \url{https://github.com/wangdongdut/DUT-Anti-UAV}
    \item BDD10K: \url{https://bdd-data.berkeley.edu/}
    \item DETRAC: \url{https://detrac-db.rit.albany.edu/}
    \item VIRAT: \url{https://viratdata.org/}
\end{itemize}

\section{Tools and Frameworks}
\begin{itemize}
    \item Ultralytics YOLO: \url{https://docs.ultralytics.com/}
    \item OpenCV: \url{https://opencv.org/}
    \item PySide6: \url{https://doc.qt.io/qtforpython/}
\end{itemize}

\section{Project Repository}
\begin{itemize}
    \item GitHub: \url{https://github.com/AEEltayeb/drone_detection_py}
\end{itemize}

\end{document}
